\section{Motivating Facts}\label{sec:motivating_facts}

Housing costs have risen sharply across advanced economies over the past
three decades.  At the same time, fertility has fallen below replacement in
nearly all of them, and the age at which women have their first child has
climbed steadily.  This section documents three stylized facts that motivate
the mechanism studied in the paper.

%% ──────────────────────────────────────────────────────────────────────────
\subsection*{Fact 1: Real house prices have risen substantially, but
heterogeneously}

Figure~\ref{fig:motivating_prices} plots real house price indices for
eight major economies, each re-indexed to 2005\,=\,100.  Several patterns
stand out.  First, the anglophone economies---the United States, the United
Kingdom, Canada, and Australia---all experienced sustained real price
appreciation, with Canada's index nearly tripling by the early 2020s.
Second, trajectories varied widely.  Japan's real house prices remained
essentially flat over this period, reflecting the long aftermath of its
1990s asset-price collapse.  Germany saw little real growth until
approximately 2010, after which prices rose rapidly before correcting in
2022--2023.  Sweden more than doubled in real terms.  Third, despite this
heterogeneity the overall direction is common: in seven of the eight
countries shown, real house prices in the 2020s are materially above their
mid-2000s level.

\begin{figure}[tp]
  \centering
  \includegraphics[width=0.92\textwidth]{motivating_house_prices.pdf}
  \caption{Real house price indices for eight major economies, indexed so
           that 2005\,=\,100.  Sources: OECD Analytical House Price
           Indicators; FRED (Case-Shiller for US; BIS residential
           property prices for Japan, Canada, Australia).}
  \label{fig:motivating_prices}
\end{figure}

%% ──────────────────────────────────────────────────────────────────────────
\subsection*{Fact 2: Fertility has fallen below replacement everywhere}

Figure~\ref{fig:motivating_fertility} shows total fertility rates (TFR) for
the same eight countries from 1960 onward.  The dashed line marks the
conventional replacement level of 2.1 children per woman.  By the 2020s
every country in the sample is below replacement.  Japan (TFR $\approx$
1.20) and Canada ($\approx$ 1.26) are among the lowest.  The United States,
long an outlier among rich countries for its relatively high fertility, has
fallen to roughly 1.6.  Even countries that maintained near-replacement
fertility into the 2000s---such as France, Australia, and Sweden---have
declined noticeably in the past decade.

The timing of the decline differs across countries.  Japan and Germany fell
below replacement in the 1970s; the anglophone countries and France held
near 2.0 until the late 2000s before declining.  But the direction is
universal.  This common trend motivates a search for explanations that
operate across institutional settings.  Housing costs, present in every
country to varying degrees, are a natural candidate.

\begin{figure}[tp]
  \centering
  \includegraphics[width=0.92\textwidth]{motivating_fertility_rates.pdf}
  \caption{Total fertility rates for eight major economies, 1960--2024.
           The dashed line marks the conventional replacement level of 2.1.
           Sources: World Bank; Eurostat.}
  \label{fig:motivating_fertility}
\end{figure}

%% ──────────────────────────────────────────────────────────────────────────
\subsection*{Fact 3: Countries with faster house price growth delayed first
births more}

Figure~\ref{fig:motivating_scatter} links the preceding two facts.  For
each of the 27 European countries with at least five years of joint data on
house prices and the mean age at first birth, we compute the cumulative real
house price growth and the total change in the mean age at first birth over
the country's observation window.  The scatter reveals a positive
association ($r = 0.50$): countries in which real house prices rose most
also experienced the largest increases in the age at which women had their
first child.  Estonia, Lithuania, and Croatia---all with cumulative real
price growth exceeding 150\%---saw first-birth ages rise by more than 3.5
years.  At the other extreme, Italy and the United Kingdom, with slower or
more moderate price growth, recorded smaller shifts.

This cross-country association is, of course, only a correlation.  It could
be driven by income growth, urbanisation, educational expansion, or other
omitted variables correlated with both housing costs and fertility timing.
We do not interpret it as causal.  Instead, we view it as a motivating
pattern that a well-specified model should be able to rationalise.  The
structural model in Section~\ref{sec:model} provides the link: when housing
is scarce and expensive, young adults delay purchasing a home, which in turn
delays---and in some cases discourages---the transition to parenthood.

\begin{figure}[tp]
  \centering
  \includegraphics[width=0.78\textwidth]{motivating_scatter_prices_timing.pdf}
  \caption{Cumulative real house price growth (\%) vs.\ change in mean age
           at first birth (years) across 27 European countries.  Turkey is
           excluded because its cumulative house-price growth exceeds 500\%,
           roughly four times the next-highest country; including Turkey
           raises $r$ to 0.54.  Each point is one country; both variables
           are computed from the first to the last year in which house prices
           and first-birth age are jointly observed.  The dashed line is a
           linear OLS fit; $r$ is the Pearson correlation.
           Sources: OECD, Eurostat.}
  \label{fig:motivating_scatter}
\end{figure}

%% ──────────────────────────────────────────────────────────────────────────

These three facts---rising real house prices, falling fertility below
replacement, and a positive cross-country association between housing cost
growth and delayed first births---motivate a model in which housing scarcity
operates as a binding constraint on family formation.  The next section
places these patterns in the related literature, and Section~\ref{sec:model}
then develops the model.
